\documentclass[11pt]{article}
\usepackage[margin=1in]{geometry}
\usepackage{booktabs}
\usepackage{graphicx}
\usepackage[hidelinks]{hyperref}
\usepackage{amsmath}
\title{FEINT: One Attacker Schema for Attacking, Hardening and Explaining\\Flow-Based Network Intrusion Detectors\\[4pt]\large Preprint draft, not submitted}
\author{Rakshit\\Vellore Institute of Technology, Chennai}
\date{October 2026}
\begin{document}
\maketitle

\begin{abstract}
Flow-based network intrusion detectors routinely report 99\,\% accuracy on public benchmarks, but
those numbers say little about an attacker who knows the model and shapes its own traffic. We
present FEINT, an open-source pipeline in which a single declarative \emph{attacker-capability
schema} (which flow features the attacker can raise, set freely, or not touch, plus the integer,
range and cross-feature relations that make a flow realisable) drives an adaptive evasion attack,
adversarial training, a robust-feature model, a poisoning sanitiser and counterfactual
explanations. Over five seeds on CIC-IDS2017 and UNSW-NB15, undefended XGBoost keeps
0.405 [0.376, 0.433] and 0.000 detection of attack flows under a realisable attack at the largest
budget; adversarial training of an XGBoost/MLP/autoencoder ensemble restores 0.793 [0.763, 0.823]
and 0.962 [0.957, 0.968]. A model restricted to the features the schema fixes is unaffected by this
attacker, but only conditionally on that schema, which we show is partly defender-favourable. On
label-corrected CIC-IDS2017 robustness numbers change materially; plain XGBoost trained on corrected
CIC-IDS2017 detects nothing on CSE-CIC-IDS2018's shared attack families (F1 0.000), while the
robust-feature model reaches F1 0.51, mostly on SSH brute force. We
report reproductions of two published baselines and of one adversarial-NIDS attack (A2PM), negative
results, and the limitations of the feature-space threat model.
\end{abstract}

\section{Introduction}
Machine-learning NIDS evaluations have well-documented pitfalls \cite{arp2022}: label errors in the
datasets \cite{engelen2021,liu2022}, base-rate blindness, and robustness claims made against
attacks that ignore what a network attacker can change \cite{pierazzi2020}. Prior work has studied
each piece separately: constrained adversarial examples for tabular and network data
\cite{simonetto2022,vitorino2022}, robustness through conserved features \cite{tong2019}, and
label sanitisation against poisoning \cite{paudice2018}. FEINT's contribution is to derive all of
these from \emph{one} schema, to measure with seeds and intervals what each derived defence buys,
and to isolate the schema's share with paired ablations for adversarial training and the
sanitiser. We make no claim of a new attack or defence algorithm.

\section{Threat model and schema}
The attacker controls one side of a flow and knows the model (white-box), or can query it for
scores or labels. Each dataset has a schema that labels every feature \emph{up} (may only
increase: duration, forward packets, forward payload bytes), \emph{free} (client TCP window,
IP TTL), \emph{fixed} (server-side and context features, destination port) or \emph{derived}
(rates and means recomputed from the others), with integer constraints and relations such as
$\text{bytes} \le \text{pkts}\cdot\text{MTU}$. A projection maps any perturbed vector back into
the realisable set; a validity check confirms it. Budgets $\epsilon$ are $L_\infty$ radii on the
controllable features in standardised $\log(1+x)$ space; derived features follow and can move
further, so the realised all-feature distance can exceed $\epsilon$.

The schema also fixes features a client can influence indirectly: UNSW-NB15 \texttt{ct\_*}
counters (by pacing or decoy connections), \texttt{trans\_depth} (HTTP pipelining) and
server-response sizes. Results for the robust-feature model are therefore conditional on the
schema.

\section{Method}
\textbf{Detectors.} XGBoost \cite{chen2016}, an MLP with analytic input gradients, a
benign-trained autoencoder, and an ensemble $\max(\text{mean}(\text{XGB},\text{MLP}),\text{AE})$;
logistic regression, random forest and isolation forest as baselines.
\textbf{Attack.} Projected gradient descent \cite{madry2018} on the target (or, for trees, on
both the undefended and the adversarially trained MLP as surrogates) followed by 100 score-based
black-box queries, keeping the per-flow best; detection at $\epsilon$ is the union over all
evaluated budgets up to $\epsilon$, so curves are monotone. Following \cite{tramer2020}, hardened
models are attacked adaptively.
\textbf{Defences.} Adversarial training with schema-projected examples; a robust-feature XGBoost
on the features no controllable or derived feature reaches; a kNN sanitiser that flags
benign-labelled training points whose neighbours in robust-feature space are malicious
(a restriction of \cite{paudice2018}); an adversarial-input alarm.
\textbf{Explanations.} Exact TreeSHAP \cite{lundberg2020} and counterfactuals defined as
minimal-budget schema-valid evasions.
\textbf{Stealing.} An MLP surrogate trained on hard-label queries \cite{tramer2016} is used for a
pure transfer attack.

\section{Setup}
CIC-IDS2017 \cite{sharafaldin2018}: 2.83\,M flows, de-duplicated in the 11-feature schema
projection, 10\,\% per class with a 5{,}000 floor (253\,k flows), stratified 70/30 split, the
subsample redrawn per seed. Corrected CIC-IDS2017 and CSE-CIC-IDS2018 from the re-release of
\cite{liu2022}/\cite{engelen2021} and the official S3 bucket (three day files, SHA-256 pinned).
UNSW-NB15 \cite{moustafa2015}: official 175\,k/82\,k partition. CICIoT2023 \cite{neto2023}: a
public ipfixprobe re-export, at most 5{,}000 flows per label for a 16\,GB machine. Five seeds,
1{,}000 held-out attack flows per curve point, 95\,\% Student-$t$ intervals ($t_4=2.78$) clipped
to $[0,1]$. All heavy runs execute on GitHub Actions runners; code, results and the workflow are
public.

\section{Results}
\begin{table}[h]\centering\small
\begin{tabular}{lcccc}\toprule
& \multicolumn{2}{c}{CIC-IDS2017} & \multicolumn{2}{c}{UNSW-NB15}\\
model & $\epsilon=0$ & $\epsilon=2$ & $\epsilon=0$ & $\epsilon=2$\\\midrule
MLP & 0.991 & 0.058 [0.001, 0.116] & 0.964 & 0.064 [0.052, 0.077]\\
Random forest & 0.995 & 0.377 [0.306, 0.447] & 0.971 & 0.000 [0.000, 0.000]\\
XGBoost & 0.998 & 0.405 [0.376, 0.433] & 0.967 & 0.000 [0.000, 0.000]\\
Ensemble & 0.998 & 0.369 [0.279, 0.459] & 0.969 & 0.328 [0.262, 0.394]\\
MLP + AT & 0.986 & 0.688 [0.580, 0.797] & 0.961 & 0.808 [0.746, 0.869]\\
XGBoost + AT & 0.999 & 0.788 [0.755, 0.821] & 0.967 & 0.941 [0.908, 0.974]\\
Ensemble + AT & 0.998 & 0.793 [0.763, 0.823] & 0.969 & 0.962 [0.957, 0.968]\\
Robust features & 0.994 & 0.994 [0.993, 0.995] & 0.942 & 0.942 [0.935, 0.950]\\\bottomrule
\end{tabular}
\caption{Detection rate of attack flows under constrained adaptive evasion, 5-seed mean [95\,\% CI].}
\end{table}

\textbf{Evasion and hardening.} Table 1 shows the collapse of undefended detectors and the recovery
from adversarial training, with clean F1 within 0.002. On CIC-IDS2017 the ensemble and XGBoost are
statistically indistinguishable at $\epsilon=2$ while the ensemble has about four times the FPR.
The robust-feature model costs 2.2 FPR points on CIC-IDS2017 and 7.1 on UNSW-NB15.

\textbf{Poisoning.} With 2\,\% of a 100\,k-flow training subsample poisoned, backdoor success is
1.00 on both datasets; the robust-feature sanitiser lowers it to 0.08 [0.02, 0.14] on CIC-IDS2017
but only to 0.68 [0.62, 0.74] on UNSW-NB15, where it also discards 10.8\,\% of clean benign flows.

\textbf{Label correction.} On corrected CIC-IDS2017 (5 seeds) undefended XGBoost keeps 0.523
[0.490, 0.556] detection at $\epsilon=2$ instead of 0.405, random forest falls to 0.005, and
adversarially trained models reach 0.980--0.997: part of
the original dataset's apparent difficulty is label noise.

\textbf{Cross-dataset.} Training on corrected CIC-IDS2017 and testing on CSE-CIC-IDS2018 over shared
families gives F1 0.000 for XGBoost, 0.211 after adversarial training and 0.506 for the
robust-feature model (5 seeds); the reverse direction gives 0.015, 0.488 [0.078, 0.898] and 0.024.

\textbf{Negative result: CICIoT2023.} On the ipfixprobe re-export no model separates benign from
attack captures (benign FPR 0.82--0.997); we do not report it as a detector.

\textbf{Schema ablation.} Two paired ablations (5 seeds, shared split, poison set and evaluation
attack) isolate what the schema buys. An MLP adversarially trained on schema-projected examples keeps
0.676 (CIC) and 0.808 (UNSW) constrained detection at $\epsilon=2$; trained on textbook
unconstrained examples at the same budget it keeps 0.031 and 0.080 (paired differences +0.645
[+0.529, +0.762] and +0.728 [+0.665, +0.792]), with no FPR difference. The kNN sanitiser in
robust-feature space leaves backdoor success at 0.083 (CIC) and 0.676 (UNSW); the same sanitiser on
all features \cite{paudice2018} leaves 1.000 and 0.903 (differences $-0.917$ [$-0.976$, $-0.858$] and
$-0.227$ [$-0.311$, $-0.144$]), at a cost of about one point of clean accuracy.

\textbf{Stealing.} With 2{,}000 label queries a stolen surrogate agrees with the victim on 95.6\,\%
(CIC) and 91.4\,\% (UNSW) of test flows and transfers at least as well as a surrogate trained on
true labels at $\epsilon=2$ (single seed).

\textbf{Reproductions.} A decision tree on the official UNSW-NB15 partition with all 42 features
reaches 0.865 accuracy and 0.248 false-alarm rate versus the 0.856 and 0.158 reported for
\cite{moustafa2016} (checked against secondary sources only: the primary article is paywalled, so
we cannot confirm its table or its false-alarm definition, and the 0.09 gap in false-alarm rate is
unexplained; for our tree neither the mean of false-positive and false-negative rates (0.146), the
false-discovery rate (0.175) nor the error rate (0.135) gives 0.158). Our reproduction of the
\cite{sharafaldin2018} Table 4 setup (multi-class, union of the per-attack selected features,
stratified 70/30 split with duplicates kept, scikit-learn defaults; split and settings assumed)
gives weighted F1 0.999 for RF (paper 0.97), 0.995 KNN, 0.998 ID3, 0.888 AdaBoost, 0.981 MLP and
0.141 naive Bayes; QDA with defaults could not be fitted (singular benign covariance), and a regularised QDA (reg\_param $10^{-3}$, a deviation) gives 0.869 (paper 0.92). On corrected CIC-IDS2017 \cite{engelen2021} RF reaches 0.995. Every fitted model except the
regularised QDA has a higher F1 than published, which we attribute mainly to duplicate flows shared
across a random split.

\begin{table}[h]\centering\small
\begin{tabular}{lcccccc}\toprule
& \multicolumn{2}{c}{targeted accuracy} & \multicolumn{2}{c}{untargeted accuracy} &
\multicolumn{2}{c}{untargeted macro-F1}\\
model & paper & ours & paper & ours & paper & ours\\\midrule
RF & 0.000 & 0.000 & 0.000 & 0.000 & 0.209 & 0.234\\
RF + AT & 0.999 & 0.932 & 0.900 & 0.901 & 0.543 & 0.642\\
MLP & 0.104 & 0.404 & 0.009 & 0.078 & 0.183 & 0.166\\
MLP + AT & 0.944 & 0.947 & 0.789 & 0.931 & 0.510 & 0.603\\\bottomrule
\end{tabular}
\caption{A2PM reproduction: scores after 50 attack iterations on CIC-IDS2017 Tuesday--Wednesday
(seed 0). Accuracy is on malicious evaluation flows, macro-F1 on all evaluation flows; ``paper'' are
the end points printed in Figures 5--7 of \cite{vitorino2022}.}
\end{table}

\textbf{Adversarial-NIDS reproduction.} We reproduce the Enterprise scenario of A2PM
\cite{vitorino2022}: CIC-IDS2017 Tuesday and Wednesday (1.14\,M flows, 8 classes), a stratified
70/30 split, the paper's random forest and Keras MLP with regular and A2PM adversarial training (one
A2PM example per malicious training flow), and 50-iteration targeted (towards benign) and untargeted
A2PM attacks adapted to the evaluation set. Table 2 shows the outcome. Regularly trained models
collapse and A2PM adversarial training preserves most of the accuracy, as in the paper. Six of the
twelve end points agree within 0.03. The largest differences are in the MLPs: the regular MLP resists
the targeted attack longer (0.404 against 0.104, although its first-iteration drop of 16 points
matches the paper's 15), and the adversarially trained MLP keeps 0.931 under the untargeted attack
against 0.789. Our adversarially trained models also keep a higher untargeted macro-F1, our
adversarially trained random forest loses 0.07 under the targeted attack where the paper's loses
nothing, and our MLP's clean macro-F1 is 0.88 against the paper's 0.97. The deviations we know of: the
MachineLearningCVE CSVs carry the destination port as one numeric column without a protocol
feature, while the paper one-hot encoded both; the MLP's input scaling and early-stopping patience
are not reported, so they are our choices. a2pm 1.2.0 perturbs one value per Python call. We
therefore use vectorised pattern subclasses that draw from the same distributions and check them
against the library in every run (largest end-to-end accuracy difference 0.0017 over four seeds).
An independent run with the library's own transforms gives the same random-forest results. A2PM
also perturbs server-side features that FEINT's schema treats as fixed, so its numbers and FEINT's
are not comparable.

\section{Limitations}
Feature-space only: no packets are rewritten and no flow extractor is re-run
\cite{pierazzi2020}. The constrained and textbook attacks are not budget-matched. The alarm and
counterfactual figures use a single budget, come from a single seed (seed 0), as do the
model-stealing figures and the A2PM reproduction, and are optimistic. UNSW-NB15 derived features are
recomputed with a convention that differs from the published columns. All robustness numbers are
empirical upper bounds from our attacks.

\section{Ethics and reproducibility}
FEINT perturbs feature vectors of public datasets against locally trained models; it generates no
traffic and contains no exploit code. Code, the benchmark workflow, per-seed result files and
pinned dataset hashes are at \url{https://github.com/rakshit-737/feint}.

\section*{Acknowledgements}
Code, experiments and this draft were produced with substantial assistance from an AI system
(Claude, Anthropic); the author is responsible for all content.

\bibliographystyle{plain}
\bibliography{refs}
\end{document}
